\begin{lemma}[The Laplacian of a polynomial function]
Let $n\in\mathbb{N}$, let $P\in\mathbb{R}[X_1,\dots,X_n]$, let $\mathbb{R}^n$
carry the Euclidean inner product, and let $u:\mathbb{R}^n\to\mathbb{R}$ be the
polynomial function $u(y)=P(y_1,\dots,y_n)$. Then for every $x\in\mathbb{R}^n$
\[
  \Delta u(x)=\sum_{i=1}^{n}\bigl(\partial_i\partial_iP\bigr)(x).
\]
\end{lemma}
\begin{proof}
Let $e_1,\dots,e_n$ be the standard orthonormal basis of $\mathbb{R}^n$, so that
$(e_k)_i=1$ when $i=k$ and $(e_k)_i=0$ otherwise. By the background proposition
computing the Laplacian from an orthonormal basis, applied with this basis,
\begin{equation}\label{eq:mpl-basis}
  \Delta u(x)=\sum_{k=1}^{n}D^2u(x)[e_k,e_k].
\end{equation}

Fix $1\le k\le n$. By \noderef{MvPolyHessian} applied with $v=w=e_k$,
\[
  D^2u(x)[e_k,e_k]
  =\sum_{i=1}^{n}\sum_{j=1}^{n}(e_k)_i\,(e_k)_j\,(\partial_i\partial_jP)(x).
\]
In this double sum the factor $(e_k)_i$ vanishes unless $i=k$, and the factor
$(e_k)_j$ vanishes unless $j=k$; so every term with $(i,j)\ne(k,k)$ is zero, and
the single surviving term, at $i=j=k$, equals
$1\cdot1\cdot(\partial_k\partial_kP)(x)$. Hence
\begin{equation}\label{eq:mpl-diag}
  D^2u(x)[e_k,e_k]=(\partial_k\partial_kP)(x)
  \qquad(1\le k\le n).
\end{equation}

Substituting \eqref{eq:mpl-diag} into \eqref{eq:mpl-basis} gives
\[
  \Delta u(x)=\sum_{k=1}^{n}(\partial_k\partial_kP)(x),
\]
which is the assertion (the summation index being a dummy).
\end{proof}
